\chapter{結論}
本研究では，YouTubeライブチャットコメントを新たなSNSコメント収集ソースとして提案し，災害時の市民感情把握への応用を目指した分析基盤の構築特徴を示した．

提案した収集・分析基盤は，低コストかつスケーラブルな構成により，継続的なデータ蓄積と多角的な分析を実現した．
約9ヶ月間の運用により，システムの安定稼働と十分なデータ収集が可能であることを確認した．

収集したデータの特性分析を通じて，短文による即応的な反応，時間的・ユーザ特性の定量化，地震発生後の反応パターンなど，重要な知見が得られた．
本研究により，YouTubeライブ配信コメントは，従来のSNSデータとは異なる特性を持つ新たなヒューマンセンサデータソースとして，災害情報収集において一定の有効性を持つことが示された．

これらの成果は，災害時における市民の情動変化をリアルタイムに把握し，公式な災害情報を補完する新たな情報源として，YouTubeライブ配信コメントが機能し得ることを示している．

今後は，リアルタイム異常検知アルゴリズムの開発や，他のデータソースとの統合により，実用的な災害検知システムへの発展が期待される．    